\section{Conclusion}\label{sec:conclusion}

We presented \method{}, a framework that unifies feed-forward generation
and per-video optimization for monocular drivable 3DGS avatars through a
continual-learning formulation. A warm-start produces a drivable avatar
within seconds of the first frame; as the monocular video streams in,
\method{} continuously refines the avatar while mitigating catastrophic
forgetting via an EWC prior on Gaussian parameters, a viewpoint-stratified
replay buffer, and pose-conditioned selective reactivation. Experiments
show that \method{} approaches the fidelity of offline optimizers, is
available almost instantly, and preserves earlier-viewpoint quality far
better than naive sequential fine-tuning.

\textbf{Limitations and future work.} SMPL registration from monocular
video remains a source of error for loose clothing and fast motion; a
learned, uncertainty-aware registration could further improve robustness.
Extending the continual-learning formulation to SMPL-X with facial and
hand articulation, and to multi-subject scenes, are natural next steps.
We also see promise in adapting more advanced continual-learning
algorithms---such as gradient-projection methods---to the 3DGS setting.

\section*{Ethical Considerations}
\method{} reconstructs photorealistic human avatars from monocular video.
Such technology carries risks of misuse including the creation of
non-consensual synthetic content. We commit to releasing the method with
documentation of these risks and support for provenance/watermarking of
generated avatars, and we encourage downstream applications to obtain
explicit subject consent.
